% An example box-and-arrow diagram, showing the house style.
% Styles (tmbox, tmboxfill, tmarrow, tmdiagram*) come from style.sty.
%
% Two rules worth keeping:
%   - tmbox (gray outline) is the default for everything;
%   - tmboxfill (pale blue) marks the ONE thing the diagram is about.
%     At most one kind of fill per diagram, or the emphasis stops meaning
%     anything.
%
% Width here stands for capacity, so the stack reads as a hierarchy
% without needing numbers on it.
\begin{tikzpicture}[x=1cm, y=1cm]

  % ---- the three tiers, smallest at the top --------------------------
  \node[tmbox, minimum width=3.8cm, minimum height=0.95cm] (top)
    at (-0.5,1.6)
    {Top tier\\{\scriptsize\color{tmgray}fast, and there is little of it}};
  \node[tmbox, minimum width=5.2cm, minimum height=0.95cm] (mid)
    at (-0.5,0.5)
    {Middle tier\\{\scriptsize\color{tmgray}slower, and there is more of it}};
  \node[tmbox, minimum width=6.6cm, minimum height=0.95cm] (bottom)
    at (-0.5,-0.6)
    {Bottom tier\\{\scriptsize\color{tmgray}slowest, and effectively unbounded}};

  % ---- who owns what: the filled box is the subject of the slide -----
  \node[tmbox, minimum width=2.6cm, minimum height=0.95cm]
    at (4.2,1.6) {Someone else};
  \node[tmboxfill, minimum width=2.6cm, minimum height=2.05cm]
    at (4.2,-0.05) {Your system};

  % ---- the axis ------------------------------------------------------
  \node[tmdiagramnote, anchor=east, text width=1.9cm] (fast) at (-4.1,1.6)
    {fastest,\\least of it};
  \node[tmdiagramnote, anchor=east, text width=1.9cm] (slow) at (-4.1,-0.6)
    {slowest,\\most of it};
  \draw[tmarrow] ([xshift=-0.3cm]slow.north east) -- ([xshift=-0.3cm]fast.south east);

\end{tikzpicture}
